\documentclass[11pt]{article}

% --- Math ---
\usepackage{amsmath, amssymb, amsthm, amsfonts}
\usepackage{mathrsfs}

% --- Encoding and units ---
\usepackage[utf8]{inputenc}
\usepackage[T1]{fontenc}
\usepackage{siunitx}

% --- Algorithms ---
\usepackage{algorithm}
\usepackage{algpseudocode}

% --- Tables ---
\usepackage{tabularx}
\usepackage{array}
\usepackage{booktabs}
\usepackage{multirow}
\usepackage{makecell}
\usepackage{longtable}
\usepackage{ragged2e}
\newcolumntype{Y}{>{\RaggedRight\arraybackslash}X}

% --- Figures ---
\usepackage{graphicx}
\usepackage{float}

% --- Bibliography ---
\usepackage[authoryear]{natbib}

\pagestyle{plain}
\setlength{\textwidth}{6.5in}
\setlength{\oddsidemargin}{0in}
\setlength{\evensidemargin}{0in}
\setlength{\textheight}{9.2in}
\setlength{\topmargin}{-0.2in}
\setlength{\headheight}{0in}
\setlength{\headsep}{0in}
\setlength{\footskip}{.4in}
\renewcommand{\baselinestretch}{1.5}

\newtheorem{theorem}{Theorem}
\newtheorem{obs}{Observation}
\newtheorem{proposition}{Proposition}
\newtheorem{remark}{Remark}
\newtheorem{definition}{Definition}

\begin{document}

\begin{center}
{\Large \bf Integrated Sample Size Optimization and Test Chamber Scheduling under Environmental Constraints}\\
\vskip 12pt
\normalsize
Nur Per\c{c}in$^{1}$,
Ali Ekici$^{1}$\\[6pt]
\footnotesize
$^{1}$Department of Industrial Engineering, Ozyegin University, Istanbul, Turkiye\\[3pt]
\texttt{nur.percin@ozu.edu.tr, ali.ekici@ozyegin.edu.tr}
\end{center}

\begin{abstract}
Refrigerator R\&D testing requires the coordinated scheduling of multiple products on heterogeneous test chambers under environmental, voltage, stage-based precedence, and due-date constraints. In current practice, sample sizes are often fixed uniformly across projects, although they directly determine pull-down workload and downstream parallelization opportunities. This study introduces the Test Chamber Scheduling with Sample Optimization Problem (TCSSOP), in which project-level sample sizes, chamber-environment assignments, and feasible schedules are optimized jointly.

An iterative framework is proposed. For a given sample vector, chamber environments are assigned according to workload distribution and chamber capabilities. Given chamber settings, sample sizes are improved through bounded local search and evaluated by an Earliest Due Date (EDD) based constructive scheduler with explicit feasibility validation. When local improvement stagnates, a Variable Neighborhood Search (VNS) diversification mechanism perturbs sample-size decisions to explore alternative workload structures.

Computational experiments on 90 benchmark scenarios compare the proposed method with fixed uniform sample policies and with variants that remove chamber-assignment feedback or VNS. The destroy ratio is selected by comparing \(0.10\), \(0.20\), and \(0.30\) at \(K_{\max}=100\). The selected setting is a single ratio of \(0.20\). It attains the lowest total tardiness in all 90 instances. Relative to the industrial three-sample policy, the average deviation is \(360.32\%\). A uniform six-sample policy reduces tardiness but remains substantially inferior, with an average deviation of \(196.95\%\). Removing feedback or VNS leaves average deviations of \(139.16\%\) and \(177.67\%\). The results show that sample size should be treated as a tactical planning decision and coordinated with chamber reconfiguration rather than fixed uniformly across projects.
\end{abstract}

\textbf{Keywords:} Test chamber scheduling; Sample size optimization; Decision-dependent workload; Resource-constrained scheduling; Variable neighborhood search; Multi-stage testing systems

\section{Introduction}
\label{sec:introduction}

The refrigerator industry operates in a highly competitive global market characterized by compressed innovation cycles and increasingly stringent market-specific requirements for performance, safety, and energy efficiency \citep{IEA2021,IEC60335-1,IEC60335-2-24,IEC62552}. Since products cannot proceed to certification, production, or commercialization until design approval tests have been successfully completed, research \& development (R\&D) testing has become a critical operational bottleneck. Delays in testing postpone market entry and reduce firms' responsiveness to changing market conditions, while inefficient utilization of capital-intensive testing infrastructure diminishes the return on R\&D investment \citep{OECD_MSTI2024,OECD_MSTI2025}. Furthermore, launch timing and quality assurance are closely linked, as the trade-off between accelerated time-to-market and product quality has a direct impact on commercial performance \citep{KimKoenigsbergOfek2022}. This study is motivated by refrigerator R\&D testing operations at a white goods manufacturer located in Turkey.

Refrigerator R\&D testing comprises a multi-stage portfolio of laboratory tests performed under controlled environmental conditions. The catalog includes gas amount determination, pull-down, energy-efficiency measurement, performance evaluation at climatic-class temperatures, optional freezing-capacity and temperature-rise tests, and consumer-usage tests. Because refrigerators must demonstrate compliance under the ambient conditions of their target markets, tests are executed at different temperature set-points (commonly spanning \(10^\circ\mathrm{C}\) to \(43^\circ\mathrm{C}\)), under ambient or humidity-controlled conditions (including \(85\%\) relative humidity for selected consumer-usage tests), and, when required by regional electrical standards, under voltage-compatible chamber settings. These changing conditions exist because climatic class, energy regulation, and electrical specification differ across regions; a single ambient environment cannot certify a product for all markets. Individual tests are long relative to daily planning cycles: typical durations range from about 5 days (pull-down and temperature rise) to 10--15 days (gas determination, performance, energy efficiency, and consumer usage). Consequently, each assignment occupies scarce chamber capacity for an extended period, and sequencing mistakes cannot be corrected quickly. The detailed test list and durations are presented in Section~\ref{sec:problem_definition}.

Which tests are required for a given product, and when they must be completed, are determined by commercial and regulatory inputs rather than by a fixed universal checklist. Target-market region and climatic class select the applicable performance and energy-efficiency variants; electrical specifications determine whether voltage-compatible execution is mandatory; product engineers may add optional tests; and project due dates reflect planned launch windows and certification milestones. Scheduling must therefore coordinate heterogeneous, product-specific test mixes under shared due-date pressure. Sample-size determination is an integral part of this decision process. In current practice, a uniform three-sample policy is often used because it is simple to administer. With few samples, upstream pull-down workload remains limited, but downstream tests that could otherwise run in parallel on different physical prototypes are serialized, which can increase completion times for due-date-critical projects. With several samples, downstream parallelism improves and tardiness may decrease, yet each additional sample multiplies pull-down operations and intensifies congestion in temperature-constrained environments. The net effect is therefore non-monotone: increasing samples does not guarantee better on-time performance, and the value of an extra sample depends on due dates, required test mixes, chamber eligibility, and stage-based release logic.

The supporting testing infrastructure consists of heterogeneous climate chambers, each containing multiple stations that operate under a shared chamber-level environment. Chambers differ in station capacity and in their technical ability to adjust temperature, humidity, and voltage. Temperature control is generally available, whereas humidity and voltage adjustment are restricted to selected chambers; thus, not every test can be executed in every chamber. Because all stations in a chamber share one environment, chamber configuration and station-level scheduling are tightly coupled. Although chambers can be reconfigured, frequent environment changes are undesirable: temperature and humidity transitions require stabilization time, interrupt ongoing tests that need a stable ambient setting, and reduce effective capacity. In practice, chamber environments are therefore treated as planning decisions that should remain stable over the scheduling horizon rather than as freely changing machine states. Chamber capabilities and the chamber--station structure are detailed in Section~\ref{sec:problem_definition}.

The operational problem we seek to solve is the following. Given a portfolio of refrigerator development projects with product-specific test requirements and due dates, decide (i)~how many physical samples to allocate to each project, (ii)~which environmental configuration to assign to each chamber, and (iii)~how to schedule all resulting test operations on chamber stations so that environmental, humidity, voltage, station-capacity, sample-availability, and stage-based precedence constraints are respected. The primary goal is to minimize total tardiness; among solutions with equal total tardiness, lower total sample usage is preferred. We refer to this integrated planning problem as the \emph{Test Chamber Scheduling with Sample Optimization Problem} (TCSSOP). Unlike classical scheduling problems with a fixed job set, sample decisions endogenously generate part of the workload and simultaneously create or restrict parallelization opportunities.

TCSSOP is challenging because its main decisions are interdependent and the induced workload is decision-dependent. Sample sizes change both the volume of pull-down work and the feasible degree of downstream parallelism; chamber-environment assignments determine which workloads can be processed where; and stage-based precedence means that early congestion propagates into later stages. Heterogeneous chamber capabilities create eligibility bottlenecks for humidity- and voltage-restricted tests, while long test durations amplify the cost of poor allocations. From an operations research perspective, the setting combines resource-constrained multi-project scheduling, unrelated parallel resources with eligibility restrictions, and tardiness objectives---features that are already computationally demanding and typically require specialized heuristic methods \citep{HartmannBriskorn2022,GomezSanchezLallaRuizFernandezGilCastroVoss2023,UlagaDurasevicJakobovic2022,MaeckerShenMonch2023,AfzaliradRezaeian2016}. The closest related problems arise in industrial test-laboratory scheduling, automotive prototype-test planning, and thermal-chamber scheduling, reviewed in Section~\ref{sec:literature_review}. Those studies still treat the testing workload as given, even when they jointly decide how many physical units are available. Reliability demonstration testing optimizes sample size under statistical criteria, but without chamber-level scheduling feasibility \citep{Fernandez2022}.

This study makes four main contributions. First, it benchmarks fixed uniform sample policies against endogenous sample-size optimization and quantifies their effect on total tardiness. Second, it includes an equal-budget control experiment in which a uniform six-sample policy uses the same average prototype consumption as the optimized solution, thereby isolating the effect of project-level allocation from total sample quantity. Third, it develops an iterative framework that coordinates sample-size decisions, chamber-environment assignment, and feasible scheduling. Fourth, it introduces a Variable Neighborhood Search (VNS) diversification strategy that perturbs workload-defining sample-size decisions rather than only sequencing or assignment decisions.

To solve TCSSOP, we propose an iterative framework. For a given sample vector, chamber environments are first generated through workload-balancing logic and further refined by local search. Based on these resulting chamber settings, sample sizes are improved through bounded local search, with each candidate evaluated by an Earliest Due Date (EDD) based constructive scheduler that explicitly validates feasibility. Since sample-size decisions change the workload distribution across environments, chamber assignments are re-optimized iteratively. When local improvement stagnates, a VNS diversification mechanism perturbs sample sizes and re-optimizes the resulting workload structure.

The remainder of the paper is organized as follows. Section~\ref{sec:literature_review} reviews related literature. Section~\ref{sec:problem_definition} defines the problem and notation, including the test catalog and chamber features. Section~\ref{sec:solution_methodology} presents the solution methodology. Section~\ref{sec:computational_study} reports computational results, and Section~\ref{sec:conclusion} concludes the study.

\section{Literature Review}
\label{sec:literature_review}

% Overleaf compile note: in running text write R\&D, never R&D.
% An unescaped & is an alignment tab and raises "Misplaced alignment tab character &".

Industrial scheduling problems have been widely studied because of their
direct impact on resource utilization, delivery performance, and operational
efficiency.
Practical settings often involve heterogeneous resources, due-date pressure,
precedence relations, eligibility restrictions, and setup-related effects,
which make exact optimization difficult and motivate specialized heuristic
and metaheuristic approaches.
These challenges become particularly pronounced in industrial testing
environments, where multiple projects compete for specialized test resources
and feasibility depends not only on resource capacity but also on technical
compatibility, test dependencies, and the availability of physical test units.
These features are central to refrigerator {R\&D} testing, in which climatic
chambers, physical samples, and multi-stage test portfolios must be
coordinated.
The closest literature can be grouped into four streams: scheduling in
industrial test laboratories and thermal-chamber facilities; the allocation
of expensive physical prototypes in {R\&D} testing; eligibility-restricted
unrelated parallel-machine scheduling; and statistical reliability-test
planning in which sample size is itself a decision.

A particularly relevant stream of research concerns scheduling in industrial
test laboratories and product-testing facilities.
\citet{MischekMusliu2021} study a real-world Test Laboratory Scheduling
Problem (TLSP) in which testing activities require heterogeneous equipment
and qualified personnel and are subject to temporal and resource-related
constraints.
An important characteristic of this setting is that individual tasks are
grouped into jobs before detailed scheduling.
Building on this structure, \citet{DanzingerEtAl2020} integrate variable
task grouping with scheduling, while \citet{MischekMusliuSchaerf2023}
further investigate local-search approaches for the combined problem.
These studies demonstrate that test-laboratory scheduling can require
decisions beyond conventional assignment and sequencing, particularly when
heterogeneous resources and the organization of testing activities interact.
Related applications emphasize additional features of physical testing.
\citet{MartinEtAl2022}, for example, consider the scheduling of a
predetermined series of product tests required for products intended to
operate under extreme environmental conditions, incorporating
technician-related constraints.
\citet{LittleEtAl2006} schedule temperature-cycling tests of electronic
circuit packs in thermal chambers, constructing feasible pack configurations
and chamber runs under shared thermal conditions.
Oven and thermal-chamber studies similarly construct feasible batches under
a shared processing environment: jobs may occupy the same resource only if
their requirements are compatible
\citep{LacknerMrkvickaMusliuWalkiewiczWinter2023,MajumderLahaSuganthan2018}.
Such applications are closely related to climatic-chamber operations because
test feasibility depends on specialized testing resources and test-specific
operating conditions.
In these studies the set of test types, and therefore the number of
operations, is given before scheduling.
In the present problem the required test types of each product are also
known, but the number of operations is not: pull-down is executed once per
physical sample, so the sample-size decision changes both the upstream
workload and the feasible downstream parallelism.

Another closely related research stream considers {R\&D} testing in which
expensive physical prototypes must be allocated among multiple tests.
\citet{BartelsZimmermann2009} study test scheduling in automotive {R\&D}
projects, accounting for temporal dependencies and resource restrictions
while seeking to reduce the number of experimental vehicles required.
\citet{LimtanyakulSchwiegelshohn2012} similarly address the assignment and
scheduling of dependent tests on vehicle prototypes using
constraint-programming and hybrid approaches.
\citet{ShiReichEpelmanKlampflCohn2017} examine prototype-vehicle test
scheduling in an automotive product-development setting, where
test--vehicle compatibility, timing restrictions, and the high cost of
prototypes motivate the coordinated allocation of tests to physical vehicles.
More recently, \citet{ClautiauxEssodaiguiNguyenSadykovYounes2025} jointly
consider the number and configurations of prototype vehicles together with
test assignment and sequencing.
Their model lexicographically minimizes the number of late tests and then
the number of prototypes used, thereby explicitly linking the planning of
physical test resources with subsequent test execution.
These studies establish that the number, characteristics, and allocation of
physical test units can have a substantial effect on the resulting {R\&D}
test schedule.

Climatic-chamber testing additionally requires the coordination of
heterogeneous test resources whose technical capabilities and operating
environments determine test eligibility.
This feature relates to scheduling models with eligibility and resource
restrictions.
For example, \citet{AfzaliradRezaeian2016} consider unrelated
parallel-machine scheduling with resource requirements, precedence
constraints, sequence-dependent setups, and machine-eligibility
restrictions, while \citet{MaeckerShenMonch2023} study an
eligibility-constrained unrelated parallel-machine problem with a
tardiness-oriented objective.
These models provide relevant scheduling foundations for heterogeneous
resources on which only subsets of operations are feasible.
In climatic-chamber testing, however, eligibility also depends on a shared
chamber-level environmental configuration: multiple stations within the same
chamber operate under a common temperature and humidity setting, while
chamber capabilities may further restrict humidity- or voltage-sensitive
tests.
Chamber configuration therefore determines the effective processing capacity
available for different portions of the testing workload.
Although eligibility-dependent scheduling models capture parts of this
interaction, they generally assume that the set of operations to be
processed is known in advance.

A separate test-planning literature treats the number of physical test units
as an endogenous decision.
In reliability demonstration testing, \citet{Fernandez2022} jointly
optimizes sample size, test duration, and acceptance criteria under
statistical risk and cost considerations.
Such models demonstrate that sample size need not be treated as a fixed
input to test planning.
Their focus, however, is primarily on the statistical design of the test
rather than on the detailed execution of multiple projects on heterogeneous
and capacity-constrained testing resources.
Thus, reliability test-planning models endogenize sample quantity but
abstract from detailed resource scheduling, whereas industrial laboratory
and prototype-test scheduling models represent operational execution in
greater detail but take both the test types and the number of operations
as given.

The problem considered in this study connects these research streams through
a different role of the sample-size decision.
For each refrigerator development project, an additional physical sample
does not merely provide another unit on which an existing set of tests can
be performed.
Because the pull-down test must be performed for every physical sample,
increasing the sample count also generates additional upstream testing
workload.
At the same time, the additional samples can be allocated to different
downstream tests, increasing the opportunities for parallel execution.
Conversely, reducing the sample count decreases pull-down workload but can
restrict downstream parallelism.
The effect of sample size on scheduling performance is therefore
non-monotone and depends on project due dates, test requirements,
physical-sample availability, stage-based precedence relations, and the
distribution of capacity across compatible chamber environments.

This mechanism distinguishes the present setting from the individual
research streams discussed above.
Prototype-test scheduling models account for limited and costly physical
test units and, in some cases, jointly determine their number or
configuration, but the set of test types and the number of operations
remain given.
Industrial test-laboratory and thermal-test scheduling models capture
heterogeneous resources, compatibility restrictions, and detailed execution
constraints, but likewise treat both the test types and the number of
operations as given.
Reliability-test planning explicitly optimizes sample size, but does not
represent the resulting detailed chamber-level execution.
In the present problem, by contrast, sample allocation simultaneously
changes the amount of upstream work and the degree of downstream
parallelism, while chamber-environment assignments determine where the
resulting workload can be processed.
The reviewed literature therefore does not jointly capture this combination
of decision-dependent workload generation, shared chamber-environment
configuration, physical-sample availability, and detailed test scheduling
under due dates.
This motivates an integrated formulation in which project-specific sample
allocation, chamber configuration, and operational scheduling are
coordinated rather than treated as independent planning decisions.

\section{Problem Definition and Notation}
\label{sec:problem_definition}

\subsection{Industrial problem setting}
\label{subsec:industrial_problem_setting}

This study is motivated by a refrigerator R\&D testing environment in which newly designed products must complete standardized laboratory tests before market release. Tests are conducted in a facility composed of climate-controlled chambers, each containing multiple parallel stations. The set of chambers is denoted by \(\mathcal{C}\). Each chamber \(c\in\mathcal{C}\) contains a fixed station set \(\mathcal{K}_c\), and all stations in the same chamber operate under a common chamber-level environmental setting.

A set of products \(\mathcal{P}\) must be processed. Each product \(p\in\mathcal{P}\) represents a development project with due date \(d_p\). The completion time \(C_p\) is defined as the finish time of the last required test of product \(p\). Required tests are product-dependent: each product has a subset \(\mathcal{T}_p\subseteq\mathcal{T}\), determined by climatic class, target-market energy regulations, electrical specifications, and optional engineering decisions.

Test execution is constrained by chamber compatibility. Each test requires a specific temperature-humidity environment, and some products require voltage-compatible execution. Since a chamber can operate under only one environmental configuration at a time and chamber capabilities are heterogeneous, not every test can be processed in every chamber. Humidity-controlled tests require humidity-capable chambers, and voltage-sensitive products can only be assigned to voltage-capable chambers.

The testing process follows multi-stage stage-based precedence logic. Gas amount determination precedes pull-down testing, pull-down acts as a release condition for downstream tests, and final consumer-usage tests are released according to downstream-stage conditions. Thus, readiness is governed by product-stage states rather than by a simple fixed predecessor list.

A central feature of the problem is the use of physical samples. Each product is represented by a sample count \(s_p\). Pull-down tests are executed once per sample, whereas most downstream tests are executed once per product and can be assigned to available samples after release. Therefore, increasing \(s_p\) raises upstream workload but may improve downstream parallelism; reducing \(s_p\) decreases workload but may restrict parallel execution. In industrial practice, sample counts are often set uniformly across projects because such rules are easy to administer. This uniform policy ignores heterogeneity in due dates, required test mixes, and chamber eligibility. Moreover, two sample vectors with the same total consumption may induce very different workload profiles when pull-down operations are concentrated on late or environment-constrained projects. This decision-dependent workload structure differentiates TCSSOP from classical scheduling problems with fixed job sets.

Overall, TCSSOP combines three coupled elements: heterogeneous chamber resources, product-dependent test requirements, and sample-size-dependent workload generation. The planner must jointly determine chamber configurations, project-level sample sizes, and feasible schedules under environmental, voltage, station-capacity, sample-availability, stage-based precedence, and due-date constraints.

\subsection{Test requirements driven by climatic class and target market}
\label{subsec:test_requirements}

Test requirements are product-dependent and are determined by the target market, climatic class, energy regulation, electrical specification, and optional engineering decisions. Therefore, each product may require a different subset of tests and a different set of chamber environments. Table~\ref{tab:test_attributes} summarizes the industrial test catalog, including stage membership, environmental conditions, sample usage, and typical durations.

\begin{table}[H]
\caption{Tests and their required environmental conditions, number of samples, and durations}
\label{tab:test_attributes}
\centering
\footnotesize
\setlength{\tabcolsep}{4pt}
\begin{tabular}{@{}cclcccc@{}}
\toprule
\textbf{Stage} &
\shortstack{\textbf{Test}\\\textbf{No.}} &
\textbf{Test Name} &
\textbf{Temperature} &
\textbf{Humidity} &
\textbf{Samples} &
\shortstack{\textbf{Duration}\\\textbf{(days)}} \\
\midrule
1 & 1 & \shortstack[l]{Refrigerant Charge\\Determination} & \shortstack{43/38/32\\\si{\degreeCelsius}} & Ambient & 1 & 10 \\
2 & 2 & Pull-down & \shortstack{43/38/32\\\si{\degreeCelsius}} & Ambient & \shortstack{All\\samples} & 5 \\
3 & 3$^{\dagger}$ & Energy Efficiency & \SI{16}{\degreeCelsius} & Ambient & 1 & 15 \\
3 & 4$^{\dagger}$ & Energy Efficiency & \SI{25}{\degreeCelsius} & Ambient & 1 & 15 \\
3 & 5$^{\dagger}$ & Energy Efficiency & \SI{32}{\degreeCelsius} & Ambient & 1 & 15 \\
3 & 6$^{*}$ & Performance & \SI{10}{\degreeCelsius} & Ambient & 1 & 10 \\
3 & 7$^{*}$ & Performance & \SI{16}{\degreeCelsius} & Ambient & 1 & 10 \\
3 & 8 & Performance & \SI{25}{\degreeCelsius} & Ambient & 1 & 10 \\
3 & 9$^{**}$ & Performance & \SI{32}{\degreeCelsius} & Ambient & 1 & 10 \\
3 & 10$^{**}$ & Performance & \SI{38}{\degreeCelsius} & Ambient & 1 & 10 \\
3 & 11$^{**}$ & Performance & \SI{43}{\degreeCelsius} & Ambient & 1 & 10 \\
3 & 12$^{***}$ & Freezing Capacity & \SI{25}{\degreeCelsius} & Ambient & 1 & 10 \\
3 & 13$^{***}$ & Temperature Rise & \SI{25}{\degreeCelsius} & Ambient & 1 & 5 \\
4 & 14 & Consumer Usage & \SI{10}{\degreeCelsius} & Ambient & 1 & 15 \\
4 & 15 & Consumer Usage & \SI{25}{\degreeCelsius} & 85\% RH & 1 & 15 \\
4 & 16 & Consumer Usage & \SI{32}{\degreeCelsius} & 85\% RH & 1 & 15 \\
\bottomrule
\end{tabular}

\vspace{1mm}
\begin{minipage}{0.95\textwidth}
\footnotesize
$^{\dagger}$ At least one of these tests is required. The energy-labeling regulation of the product's target market selects the pair at \(16^\circ\mathrm{C}\) and \(32^\circ\mathrm{C}\), or a single test at \(25^\circ\mathrm{C}\) or at \(32^\circ\mathrm{C}\).\\
$^{*}$ Low-temperature performance test of the climatic class: \(10^\circ\mathrm{C}\) for products targeted for subnormal climates and \(16^\circ\mathrm{C}\) otherwise.\\
$^{**}$ High-temperature performance test of the climatic class: \(43^\circ\mathrm{C}\) for products targeted for tropical climates, \(38^\circ\mathrm{C}\) for subtropical climates, and \(32^\circ\mathrm{C}\) for normal and subnormal climates.\\
$^{***}$ Optional. The product engineer may request either test, both, or neither.
\end{minipage}
\end{table}

Refrigerant charge determination, pull-down, the performance test at
\(25^\circ\mathrm{C}\), and the three consumer-usage tests are required
for every product.
The temperature and humidity of each test are given in
Table~\ref{tab:test_attributes}.
Ambient humidity means that the test needs no humidity adjustment and
can be run in any chamber.
An entry of 85\% RH means that the test must be run at 85\% relative
humidity, so it can be assigned only to a chamber with humidity control.
Every test uses one sample, except pull-down, which is repeated on each
sample.
For refrigerant charge determination and pull-down, the temperature is
the upper ambient limit of the product's climatic class.
Products targeted for tropical, subtropical, and normal/subnormal
climates are tested at \(43^\circ\mathrm{C}\), \(38^\circ\mathrm{C}\),
and \(32^\circ\mathrm{C}\), respectively.
In addition to the performance test at \(25^\circ\mathrm{C}\), the
climatic class determines one low-temperature and one high-temperature
performance test, as specified in Table~\ref{tab:test_attributes}.
An additional performance test at \(16^\circ\mathrm{C}\) is required for
products targeting tropical, subtropical, and normal climates, whereas
the corresponding test for products targeting subnormal climates is
performed at \(10^\circ\mathrm{C}\).
High-temperature performance tests at \(43^\circ\mathrm{C}\),
\(38^\circ\mathrm{C}\), and \(32^\circ\mathrm{C}\) are required for
products targeting tropical, subtropical, and normal or subnormal
climates, respectively.
The applicable energy-efficiency tests are determined by the
energy-labeling regulation of the product's target market.
That regulation specifies exactly one of three mutually exclusive
options: the pair of tests at \(16^\circ\mathrm{C}\) and
\(32^\circ\mathrm{C}\), a single test at \(25^\circ\mathrm{C}\), or a
single test at \(32^\circ\mathrm{C}\).
The three energy-efficiency rows in Table~\ref{tab:test_attributes} are
therefore alternatives, not a mandatory set.
Freezing-capacity and temperature-rise tests are optional and are
included when requested by the product engineer.
Thus, each product generally has a different test portfolio even though
all products follow the same stage-based testing structure.

Electrical requirements also vary by region. Typical configurations include \(230\,\mathrm{V}\)--\(50\,\mathrm{Hz}\) for Europe and Turkey, \(120\,\mathrm{V}\)--\(60\,\mathrm{Hz}\) for the USA and Canada, and \(230\,\mathrm{V}\)--\(60\,\mathrm{Hz}\) for markets such as Saudi Arabia and Brazil. Nonstandard voltage means a target-market specification other than the laboratory supply of \(230\,\mathrm{V}\)--\(50\,\mathrm{Hz}\); for example, a product destined for the USA requires \(120\,\mathrm{V}\)--\(60\,\mathrm{Hz}\) and must be tested in a voltage-capable chamber. Voltage compatibility is modeled at the product level. If product \(p\) requires voltage-compatible execution, all operations of that product must be assigned to voltage-capable chambers; otherwise, the product can be processed in any environmentally compatible chamber.

Each test type \(t\in\mathcal{T}\) is associated with an environmental requirement \(e(t)\), specifying its temperature and humidity condition.

As a result, product-specific test definitions generate heterogeneous workload patterns across environments and voltage-capable resources. This heterogeneity affects both feasibility and scheduling efficiency, since workload may concentrate in specific temperature-humidity environments while voltage-sensitive products compete for a restricted subset of chambers. Because required test subsets differ across products, the marginal value of an additional sample also differs across projects, reinforcing the need for project-specific rather than uniform sample decisions.

\subsection{Chamber structure and capability constraints}
\label{subsec:chamber_structure}

The testing facility consists of a heterogeneous set of chambers \(\mathcal{C}\). Each chamber \(c\in\mathcal{C}\) contains a fixed station set \(\mathcal{K}_c\), with capacity \(M_c=|\mathcal{K}_c|\). The industrial system considered in this study includes 28 environmental chambers and 240 testing stations. Chamber capacities range from 6 to 12 stations. Figure~\ref{fig:chamber_structure} illustrates the chamber--station structure: each chamber consists of multiple stations operating under shared environmental conditions, while chamber capabilities may differ across the facility.

\begin{figure}[H]
\centering
\includegraphics[width=1\textwidth]{chamber_structure.png}
\caption{Schematic representation of test chambers and station allocation. Each chamber consists of multiple stations operating under shared environmental conditions, while chamber capabilities may differ.}
\label{fig:chamber_structure}
\end{figure}

All stations within a chamber operate under a common chamber-level environment. Therefore, environmental compatibility is enforced at the chamber level rather than independently for each station. Detailed chamber capacities and technical adjustment capabilities are summarized in Table~\ref{tab:chamber_capabilities}.

\begin{table}[H]
\centering
\caption{Chamber capacities and technical adjustment capabilities}
\label{tab:chamber_capabilities}
\small
\setlength{\tabcolsep}{5pt}
\renewcommand{\arraystretch}{1.12}
\begin{tabular}{lccccc}
\hline
\textbf{Chambers} &
\shortstack{\textbf{No. of}\\\textbf{Chambers}} &
\shortstack{\textbf{Stations}\\\textbf{per Chamber}} &
\shortstack{\textbf{Temperature}\\\textbf{Adjustment}} &
\shortstack{\textbf{Humidity}\\\textbf{Adjustment}} &
\shortstack{\textbf{Voltage}\\\textbf{Adjustment}} \\
\hline
T001, T006--T007 & 3 & 8  & 1 & 0 & 0 \\
T002--T005, T025 & 5 & 8  & 1 & 1 & 1 \\
T008, T026--T028 & 4 & 12 & 1 & 0 & 1 \\
T009              & 1 & 12 & 1 & 0 & 0 \\
T010--T013        & 4 & 12 & 1 & 1 & 0 \\
T014--T015        & 2 & 6  & 1 & 1 & 0 \\
T016--T023        & 8 & 6  & 1 & 0 & 1 \\
T024              & 1 & 8  & 1 & 1 & 0 \\
\hline
\end{tabular}
\end{table}

The indicators in Table~\ref{tab:chamber_capabilities} show whether a chamber group supports temperature, humidity, and voltage adjustment. Temperature adjustment is available in all chambers, whereas humidity and voltage adjustment are available only in selected chambers. Let \(h_c\in\{0,1\}\) and \(v_c\in\{0,1\}\) denote the humidity and voltage capability indicators of chamber \(c\), respectively, and let \(\mathcal{C}^v=\{c\in\mathcal{C}:v_c=1\}\) be the set of voltage-capable chambers.

Let \(\mathcal{E}\) denote the set of feasible environment configurations, where each environment corresponds to a temperature-humidity combination. For each chamber \(c\), a single environment \(X_c\in\mathcal{E}\) is selected. A test operation \(t\) can be assigned to chamber \(c\) only if the chamber environment matches the required environment:
\[
X_c=e(t).
\]
In addition, tests requiring controlled humidity, such as 85\% relative humidity, can only be assigned to chambers with \(h_c=1\). Voltage compatibility is required when the product's target-market electrical standard is nonstandard relative to the laboratory supply of \(230\,\mathrm{V}\)--\(50\,\mathrm{Hz}\). For example, a product destined for the USA requires \(120\,\mathrm{V}\)--\(60\,\mathrm{Hz}\) and must be tested in a voltage-capable chamber. Let \(\mathcal{P}^v\) denote the set of such products (equivalently, those with \(\nu_p=1\)). If \(p\in\mathcal{P}^v\), all operations of product \(p\) must be assigned to a chamber in \(\mathcal{C}^v\). Products not in \(\mathcal{P}^v\) can be processed in any environmentally compatible chamber.

Because chamber capacities and capabilities are not uniformly distributed, the effective capacity available for a specific environment may be much smaller than the total station capacity. Thus, bottlenecks may arise not only from workload concentration but also from limited humidity- or voltage-capable resources. Moreover, although chambers can be reconfigured, frequent environment changes are operationally undesirable due to stabilization delays and disruption of ongoing tests; chamber environments are therefore treated as planning decisions that remain fixed over the scheduling horizon. Feasible schedules must therefore jointly satisfy environmental compatibility, humidity feasibility, voltage feasibility, station-capacity constraints, and sample-dependent workload requirements.

\subsection{Test flow, stage logic, and workload generation}
\label{subsec:test_flow}

The test process follows a structured multi-stage flow:
\[
\text{Refrigerant Charge Determination}\rightarrow \text{Pull-down}\rightarrow \text{Other Required Tests}\rightarrow \text{Consumer Usage}.
\]
Stage progression is controlled by stage-based precedence rules rather than by a simple fixed predecessor list. Gas amount determination, when required, precedes pull-down testing. Pull-down acts as the release condition for downstream tests. Other required tests can start only after the pull-down phase is completed, while consumer-usage tests are released after pull-down completion and after all required downstream tests have been initiated.

The required test set is product-specific. For each product \(p\), the set \(\mathcal{T}_p\subseteq\mathcal{T}\) is determined by climatic class, electrical specifications, target-market energy regulations, and optional engineering decisions. Since several test groups represent alternative operating conditions, only the regulation-consistent variants are selected for each product.

A central feature of TCSSOP is decision-dependent workload generation through sample counts. Each product \(p\) has a sample-size decision \(s_p\). Pull-down tests are executed once per sample, whereas other required tests are executed once per product. The multiplicity of test \(t\) for product \(p\) is therefore
\[
m_{pt}(s_p)=
\begin{cases}
s_p, & \text{if }t\text{ is a pull-down test,}\\
1, & \text{otherwise.}
\end{cases}
\]
For environment \(e\), the induced workload is
\[
W_e(S)=
\sum_{p\in\mathcal{P}}
\sum_{\substack{t\in\mathcal{T}_p\\e(t)=e}}
m_{pt}(s_p)\tau_t .
\]

This structure creates the main trade-off of the problem. Increasing \(s_p\) increases upstream pull-down workload and may create congestion in specific environments, but it can also enable downstream tests to run in parallel on different physical samples. Reducing \(s_p\) lowers workload but may restrict downstream parallelism. Hence, the effect of sample size on tardiness is non-monotone and depends on chamber assignment, sample availability, due dates, and stage-based precedence release conditions.

Unlike classical scheduling problems with a fixed job set, the number and distribution of operations in TCSSOP are induced by both the required test set \(\mathcal{T}_p\) and the sample vector \(S=\{s_p\}_{p\in\mathcal{P}}\). This decision-dependent workload generation is the main source of complexity and motivates the integrated solution framework.

\subsection{Objective and feasibility conditions}
\label{subsec:objective_feasibility}

The objective is to determine feasible test plans that complete products on time while using limited testing resources effectively. For each product \(p\in\mathcal{P}\), let \(C_p\) denote its completion time, defined as the finish time of its last required test operation. Given due date \(d_p\), product tardiness is
\[
T_p=\max(0,C_p-d_p).
\]
The primary objective is to minimize total tardiness:
\[
\min T^{tot}=\sum_{p\in\mathcal{P}}T_p.
\]
Sample usage is considered as a secondary criterion. Let \(s_p\) be the number of samples assigned to product \(p\). Total sample usage is
\[
S^{tot}=\sum_{p\in\mathcal{P}}s_p.
\]
Solutions are evaluated lexicographically: total tardiness is minimized first, and among solutions with equal total tardiness, the solution with the smaller total sample usage is preferred.

A feasible solution must satisfy all operational, environmental, and stage-based requirements of the testing system. For each product \(p\), all tests in \(\mathcal{T}_p\) must be scheduled with the correct multiplicities: pull-down tests generate \(s_p\) operations, while other required tests generate one operation per product, and no non-required tests are scheduled. Each operation must be assigned to a chamber whose environment matches its test requirement \(e(t)\). Tests requiring controlled humidity can only be assigned to chambers with humidity capability \((h_c=1)\). If product \(p\) requires voltage-compatible execution \((\nu_p=1)\), all of its operations must be assigned to voltage-capable chambers \((v_c=1)\); otherwise, the product may be assigned to any environmentally compatible chamber.

Feasibility also requires temporal and resource consistency. Each station can process at most one operation at a time, and each physical sample can participate in at most one operation at a time. Stage-based precedence rules must be respected: downstream tests can start only after the pull-down phase is completed, and consumer-usage tests require both completion of the pull-down phase and initiation of all required downstream tests. The sample size of each product is restricted to the interval \([s_{\min},s_{\max}]\). The lower bound \(s_{\min}\) is the minimum number of prototypes permitted in the planning process. A single sample is not used: every remaining test of that product would occupy the same physical unit, so those tests cannot overlap, and the plan would have no replicate if a unit is damaged or a test must be repeated. Current industrial practice therefore plans with three prototypes per product; we set \(s_{\min}=3\) to this laboratory floor rather than to a statistical sample-size formula. The upper bound \(s_{\max}\) is the maximum number of tests of a product that can run concurrently. Using more than \(s_{\max}\) samples cannot create additional downstream parallelism and would only generate additional pull-down workload. Finally, the completion time \(C_p\) of each product is the maximum completion time among all its required operations.

These conditions define feasibility at both structural and temporal levels. Structurally, test selection, sample counts, and chamber assignments must be compatible with environmental and capability requirements. Temporally, the generated operations must admit a schedule that satisfies station capacity, sample exclusivity, and stage-based precedence constraints.

\subsection{Notation}
\label{subsec:notation}

The main notation used in the definition of RCSP-SO is summarized in Table~\ref{tab:notation_RCSP-SO}. Objective values are evaluated lexicographically by prioritizing total tardiness and, among solutions with equal total tardiness, minimizing total sample usage.

\begin{table}[H]
\caption{Sets, parameters and solution measures used in the definition of RCSP-SO}
\label{tab:notation_RCSP-SO}
\label{tab:notation_tcssop}
\centering
\small
\begin{tabular}{cl}
\toprule
\textbf{Notation} & \textbf{Description}\\
\midrule
\textbf{Sets} & \\
\midrule
\(\mathcal{P}\) & Set of products, indexed by \(p\) \\
\(\mathcal{P}^v\) & Subset of products that require nonstandard voltage (e.g., USA, \(120\,\mathrm{V}\)--\(60\,\mathrm{Hz}\)) \\
\(\mathcal{C}\) & Set of chambers, indexed by \(c\) \\
\(\mathcal{C}^h\) & Subset of chambers that have humidity-adjustment capability \\
\(\mathcal{C}^v\) & Subset of chambers that have voltage-adjustment capability \\
\(\mathcal{K}_c\) & Set of stations in chamber \(c\) \\
\(\mathcal{E}\) & Set of environmental (temperature and humidity) configurations, indexed by \(e\) \\
\(\mathcal{E}^H\) & Set of environmental configurations that require 85\% RH \\
\(\mathcal{T}\) & Set of tests, indexed by \(t\) \\
\(\mathcal{T}_p\) & Set of tests required by product \(p\) \\
\midrule
\textbf{Parameters} & \\
\midrule
\(d_p\) & Due date of product \(p\) \\
\(e(t)\) & Environmental configuration required by test \(t\) \\
\(\tau_t\) & Duration of test \(t\) \\
\(K_c\) & Number of stations in chamber \(c\) \\
\(h_c\) & 1 if chamber \(c\) can adjust humidity; 0 otherwise \\
\(v_c\) & 1 if chamber \(c\) can adjust voltage; 0 otherwise \\
\(s_{\min}\) & Lower bound on the sample size of a product \\
\(s_{\max}\) & Upper bound on the sample size of a product \\
\midrule
\textbf{Solution measures} & \\
\midrule
\(C_p\) & Completion time of product \(p\) \\
\(T_p\) & Tardiness of product \(p\), \(T_p=\max\{0,C_p-d_p\}\) \\
\(T^{tot}\) & Total tardiness, \(T^{tot}=\sum_{p\in\mathcal{P}}T_p\) \\
\bottomrule
\end{tabular}
\end{table}

\section{Solution Methodology}
\label{sec:solution_methodology}

\subsection{Overall Solution Framework}
\label{subsec:framework}
\label{subsec:overall_framework}

The method is built around a simple observation: sample sizes and
chamber settings cannot be chosen independently.
Increasing the number of samples of a product creates one additional
pull-down operation for each extra unit, so the upstream workload
grows.
At the same time, extra samples allow later tests to run in parallel,
which can reduce completion times when due dates are tight.
Which of these two effects dominates depends on how chamber capacity
is currently allocated across environments.
The search therefore alternates between assigning chambers to a given
sample vector and improving that sample vector under a fixed
assignment.
Algorithm~\ref{alg:framework} records this outer coordination;
the inner procedures are developed in the following subsections.

\begin{algorithm}[H]
\caption{Overall Solution Framework}
\label{alg:framework}
\label{alg:overall_framework}
\begin{algorithmic}[1]
\State \textbf{Input:} instance data
\State \textbf{Output:} \((\mathbf{x}_{inc},\mathbf{s}_{inc},\Pi_{inc},T^{tot}_{inc})\)
\State Set \(s_p\leftarrow s_{\min}\) for all \(p\in\mathcal{P}\); \(\mathbf{s}\leftarrow(s_p)_{p\in\mathcal{P}}\) \label{step:fw-init}
\State \((\mathbf{x}_{inc},\mathbf{s}_{inc},\Pi_{inc},T^{tot}_{inc})\leftarrow(\emptyset,\mathbf{s},\emptyset,+\infty)\)
\For{\(i\leftarrow 1\) to \(I_{\max}\)}
    \State \(\mathbf{x}\leftarrow\text{ChamberAssignment}(\mathbf{s})\) \label{step:fw-assign}
    \State \((\mathbf{s}_{cand},\Pi_{cand},T^{tot}_{cand})\leftarrow\text{SampleSizeLocalSearch}(\mathbf{x},\mathbf{s})\) \label{step:fw-ls}
    \If{\((T^{tot}_{inc},\mathbf{1}^{\mathsf{T}}\mathbf{s}_{inc})\succ_{\mathrm{lex}}(T^{tot}_{cand},\mathbf{1}^{\mathsf{T}}\mathbf{s}_{cand})\)} \label{step:fw-lex}
        \State \((\mathbf{x}_{inc},\mathbf{s}_{inc},\Pi_{inc},T^{tot}_{inc})\leftarrow(\mathbf{x},\mathbf{s}_{cand},\Pi_{cand},T^{tot}_{cand})\)
        \State \(\mathbf{s}\leftarrow\mathbf{s}_{cand}\)
    \Else
        \State \((\mathbf{s}_{cand},\Pi_{cand},T^{tot}_{cand})\leftarrow\text{VNS}(\mathbf{x}_{inc},\mathbf{s}_{inc},\Pi_{inc},T^{tot}_{inc})\) \label{step:fw-vns}
        \If{\((T^{tot}_{inc},\mathbf{1}^{\mathsf{T}}\mathbf{s}_{inc})\succ_{\mathrm{lex}}(T^{tot}_{cand},\mathbf{1}^{\mathsf{T}}\mathbf{s}_{cand})\)}
            \State \((\mathbf{s}_{inc},\Pi_{inc},T^{tot}_{inc})\leftarrow(\mathbf{s}_{cand},\Pi_{cand},T^{tot}_{cand})\)
            \State \(\mathbf{s}\leftarrow\mathbf{s}_{cand}\)
        \Else
            \State \Return \((\mathbf{x}_{inc},\mathbf{s}_{inc},\Pi_{inc},T^{tot}_{inc})\) \label{step:fw-stop}
        \EndIf
    \EndIf
\EndFor
\State \Return \((\mathbf{x}_{inc},\mathbf{s}_{inc},\Pi_{inc},T^{tot}_{inc})\)
\end{algorithmic}
\end{algorithm}

We start from the minimum feasible sample count for every product
(Step~\ref{step:fw-init}).
At each outer iteration the chambers are first assigned to the
current sample vector (Step~\ref{step:fw-assign}).
Sample sizes are then improved by local search under this assignment
(Step~\ref{step:fw-ls}).
That procedure evaluates each trial by an earliest-due-date (EDD)
schedule and returns the resulting sample vector, schedule, and total
tardiness.
A candidate is accepted if
\((T^{tot}_{inc},\mathbf{1}^{\mathsf{T}}\mathbf{s}_{inc})\succ_{\mathrm{lex}}(T^{tot}_{cand},\mathbf{1}^{\mathsf{T}}\mathbf{s}_{cand})\)
(Step~\ref{step:fw-lex}).
The search then continues from the accepted samples, so the next
iteration reassigns the chambers.
If local search does not improve the incumbent, VNS is applied to
the incumbent
\((\mathbf{x}_{inc},\mathbf{s}_{inc},\Pi_{inc},T^{tot}_{inc})\)
(Step~\ref{step:fw-vns}): it perturbs the sample vector and repeats
the same local search under the incumbent assignment.
If the same lexicographic test holds after VNS, only the incumbent
samples, schedule, and tardiness are replaced; the next iteration
reassigns the chambers to the new samples.
Otherwise the search stops (Step~\ref{step:fw-stop}).
It also stops when the iteration limit \(I_{\max}\) is reached.

\subsection{Chamber assignment under environment constraints}
\label{subsec:chamber_assignment}

In this section, we explain how the chamber configurations are determined based on the test workload. We use \(\mathbf{s}=(s_p)_{p\in\mathcal{P}}\) to represent the sample-size vector. For a fixed sample-size vector \(\mathbf{s}\), we determine the environmental configuration of each chamber solving a mixed-integer linear program. We implement this chamber configuration phase in two steps. Since sample counts affect the workload of environmental configurations, chamber assignment must be updated when \(\mathbf{s}\) changes. We use \(m_{pt}(s_p)\) to represent the number of samples on which test \(t\) of product \(p\) has to be executed. We first determine the total test workload for each environmental configuration based on voltage-adjustment requirement. More specifically, for each environmental configuration \(e \in \mathcal{E}\) we calculate the overall total test duration and total test duration with nonstandard voltage as follows:
\begin{align*}
W_e(\mathbf{s}) &=
\sum_{p\in\mathcal{P}}
\sum_{\substack{t\in\mathcal{T}_p\\e(t)=e}}
m_{pt}(s_p)\tau_t,
\\
W_e^V(\mathbf{s}) &=
\sum_{p\in\mathcal{P}^v}
\sum_{\substack{t\in\mathcal{T}_p\\e(t)=e}}
m_{pt}(s_p)\tau_t.
\end{align*}
\(W_e(\mathbf{s})\) denotes the total test duration under environmental configuration \(e\), and \(W_e^V(\mathbf{s})\) denotes the total duration of the tests requiring nonstandard voltage under environmental configuration \(e\). The workload shares of each environmental configuration, considering all tests and only those requiring nonstandard voltage, respectively, are defined as follows:
\[
\phi_e=\frac{W_e(\mathbf{s})}{\sum\limits_{e\in\mathcal{E}}W_e(\mathbf{s})},\qquad
\phi_e^V=\frac{W_e^V(\mathbf{s})}{\sum\limits_{e\in\mathcal{E}}W_e^V(\mathbf{s})}.
\]
The target capacities (in terms of number of stations required) are calculated as follows:
\[
S_e=\phi_e \sum_{c\in\mathcal{C}}K_c,\qquad
S_e^V=\phi_e^V \sum_{c\in\mathcal{C}^v}K_c.
\]
We define \(\mathcal{E}_\mathbf{s}\) and \(\mathcal{E}^V_\mathbf{s}\) as the set environmental conditions for which nonzero test durations are needed considering all tests and the ones requiring nonstandard voltage only. That is,
\[
\mathcal{E}_\mathbf{s}=\{e\in\mathcal{E}:\ W_e(\mathbf{s})>0\},
\]
\[
\mathcal{E}^V_\mathbf{s}=\{e\in\mathcal{E}:W_e^V(\mathbf{s})>0\}.
\]
Table~\ref{tab:notation_ccd} lists the additional notation of the chamber configuration model. Symbols already defined in Table~\ref{tab:notation_RCSP-SO} are omitted.

\begin{table}[H]
\caption{Additional notation used in the integer programming model for determining chamber configurations}
\label{tab:notation_ccd}
\centering
\small
\begin{tabular}{cl}
\toprule
\textbf{Notation} & \textbf{Description}\\
\midrule
\(W_e(\mathbf{s})\) & Total test duration under configuration \(e\) \\
\(W_e^V(\mathbf{s})\) & Nonstandard-voltage test duration under configuration \(e\) \\
\(\mathcal{E}_\mathbf{s}\) & Set of configurations with \(W_e(\mathbf{s})>0\) \\
\(\mathcal{E}^V_\mathbf{s}\) & Set of configurations with \(W_e^V(\mathbf{s})>0\) \\
\(\phi_e\) & Workload share of configuration \(e\) \\
\(\phi_e^V\) & Nonstandard-voltage workload share of configuration \(e\) \\
\(S_e\) & Target number of stations for configuration \(e\) \\
\(S_e^V\) & Target number of voltage-adjustable stations for configuration \(e\) \\
\midrule
\textbf{Decision Variables} & \\
\midrule
\(x_{ce}\) & 1, if chamber \(c\) is assigned to configuration \(e\); 0, otherwise \\
\(z_e\) & \shortstack[l]{Deviation between the number of assigned and\\targeted stations for configuration \(e\)} \\
\(z_e^V\) & \shortstack[l]{Deviation between the number of assigned and\\targeted voltage-adjustable stations for configuration \(e\)} \\
\bottomrule
\end{tabular}
\end{table}




Using the workload calculations explained above, we propose a mixed-integer linear program to determine the configurations of the chambers. Our goal in this model is to minimize the summation of the differences between the number of stations assigned to a configuration and the target capacity for that configuration. We define \(x_{ce}\) as a binary variable that equals 1 if chamber \(c\) is assigned to environmental configuration \(e\), and 0 otherwise. Moreover, we define \(z_e\) and \(z_e^V\) as the differences between the total number of stations assigned to environmental configuration \(e\) and the corresponding target capacities, considering all tests and only those requiring nonstandard voltage, respectively.
The proposed chamber configuration determination model is as follows:
\begin{align}
\textrm{MIP-CCD($\mathbf{s}$)}: \min\quad
& \sum_{e\in\mathcal{E}}z_e+\sum_{e\in\mathcal{E}}z_e^V
\label{eq:assign_obj}\\
\text{s.t.}\quad
& \sum_{e\in\mathcal{E}}x_{ce}=1,
&& \forall c\in\mathcal{C},
\label{eq:assign_one}\\
& \sum_{c\in\mathcal{C}}x_{ce}\ge 1,
&& \forall e\in\mathcal{E}_\mathbf{s},
\label{eq:assign_cover}\\
& \sum_{c\in\mathcal{C}^v}x_{ce}\ge 1,
&& \forall e\in\mathcal{E}^V_\mathbf{s},
\label{eq:assign_volt}\\
& x_{ce}\le h_c,
&& \forall c\in\mathcal{C},\ e\in\mathcal{E}^H,
\label{eq:assign_hum}\\
& z_e\ge \sum_{c\in\mathcal{C}}K_cx_{ce}-S_e,
&& \forall e\in\mathcal{E},
\label{eq:assign_z1}\\
& z_e\ge S_e-\sum_{c\in\mathcal{C}}K_cx_{ce},
&& \forall e\in\mathcal{E},
\label{eq:assign_z2}\\
& z_e^V\ge \sum_{c\in\mathcal{C}^v}K_cx_{ce}-S_e^V,
&& \forall e\in\mathcal{E},
\label{eq:assign_zv1}\\
& z_e^V\ge S_e^V-\sum_{c\in\mathcal{C}^v}K_cx_{ce},
&& \forall e\in\mathcal{E},
\label{eq:assign_zv2}\\
& x_{ce}\in\{0,1\},
&& \forall c \in \mathcal{C}, e \in \mathcal{E},
\label{eq:assign_bin}\\
& z_e,z_e^V\ge 0,
&& \forall e \in \mathcal{E}.
\label{eq:assign_dom}
\end{align}
In this model, objective~\eqref{eq:assign_obj} minimizes the total deviation between assigned station counts and the two target capacities. Constraint~\eqref{eq:assign_one} assigns each chamber to exactly one environmental configuration. Constraint~\eqref{eq:assign_cover} gives every configuration with a positive workload at least one chamber. Constraint~\eqref{eq:assign_volt} gives every configuration with a positive nonstandard-voltage workload at least one voltage-capable chamber. Constraint~\eqref{eq:assign_hum} forbids humidity-controlled configurations on chambers that cannot hold humidity. Constraints~\eqref{eq:assign_z1}--\eqref{eq:assign_z2} force \(z_e\) to be at least the absolute difference between the number of stations assigned to \(e\) and \(S_e\). Constraints~\eqref{eq:assign_zv1}--\eqref{eq:assign_zv2} do the same for voltage-capable stations and \(S_e^V\). The remaining constraints specify the domains of the decision variables.

After the chamber configurations are determined, as a second step we implement a local search phase to fine-tune the chamber configuration decisions. We use the scheduling procedure to form a schedule according to the given chamber configurations and the sample-size vector. The normalized workload of chamber \(c\) is \(\lambda_c=B_c/K_c\), where \(B_c\) is the total duration of the tests assigned to that chamber and \(K_c\) is its number of stations. A chamber is \emph{highly loaded} if \(\lambda_c\) is large and \emph{lightly loaded} if \(\lambda_c\) is small. Chambers are ranked from the largest \(\lambda_c\) to the smallest, and pairs are scanned in this heaviest--lightest order. We consider exchanging the configurations of a highly loaded chamber and a lightly loaded chamber with \(\lambda_c>\lambda_{c'}\) (without violating \eqref{eq:assign_one}--\eqref{eq:assign_hum}). For each feasible swap we call the scheduling procedure. The first neighbor that improves the lexicographic objective is accepted immediately; the remaining swap pairs in that pass are not evaluated. If no swap is accepted in that pass, we reassign a lightly loaded chamber to the configuration of a highly loaded chamber, again in the same ranked order, and accept the first improving feasible move. After an accepted move, \(\lambda_c\) is recomputed and the next pass starts again with swap. These two moves are implemented sequentially as long as the solution improves, or until \(B_{\mathrm{room}}\) calls of the scheduling procedure have been used. The steps are provided in Algorithm~\ref{alg:chamber_assignment}.

\begin{algorithm}[H]
\caption{\textsc{ChamberAssignment}()}
\label{alg:chamber_assignment}
\begin{algorithmic}[1]
\State \textbf{Input:} sample-size vector \(\mathbf{s}\)
\State \textbf{Output:} chamber configuration assignment \(\mathbf{x}\)
\State \(\mathbf{x}\leftarrow\Call{MIP-CCD}{\mathbf{s}}\) \label{step:ca-exact}
\State \((\Pi,T^{tot})\leftarrow\Call{Schedule}{\mathbf{x},\mathbf{s}}\) \label{step:ca-eval}
\Repeat
    \State Rank chambers from largest \(\lambda_c=B_c/K_c\) to smallest using \(\Pi\) \label{step:ca-rank}
    \State \(\mathit{improved}\leftarrow\mathrm{false}\)
    \For{pairs \((c,c')\) in heaviest--lightest order with \(\lambda_c>\lambda_{c'}\)}
        \State \(\mathbf{x}_{cand}\leftarrow\Call{Swap}{\mathbf{x},c,c'}\) \Comment{exchange the configurations of \(c\) and \(c'\)} \label{step:ca-swap}
        \If{\(\mathbf{x}_{cand}\) satisfies \eqref{eq:assign_one}--\eqref{eq:assign_hum}}
            \State \((\Pi_{cand},T^{tot}_{cand})\leftarrow\Call{Schedule}{\mathbf{x}_{cand},\mathbf{s}}\)
            \If{\((T^{tot},\mathbf{1}^{\mathsf{T}}\mathbf{s})\succ_{\mathrm{lex}}(T^{tot}_{cand},\mathbf{1}^{\mathsf{T}}\mathbf{s})\)}
                \State \((\mathbf{x},\Pi,T^{tot})\leftarrow(\mathbf{x}_{cand},\Pi_{cand},T^{tot}_{cand})\) \label{step:ca-accept}
                \State \(\mathit{improved}\leftarrow\mathrm{true}\); \textbf{break}
            \EndIf
        \EndIf
    \EndFor
    \If{\(\mathit{improved}=\mathrm{false}\)}
        \For{pairs \((c,c')\) in heaviest--lightest order with \(\lambda_c>\lambda_{c'}\)}
            \State \(\mathbf{x}_{cand}\leftarrow\Call{Move}{\mathbf{x},c,c'}\) \Comment{assign \(c'\) to the configuration of \(c\)} \label{step:ca-move}
            \If{\(\mathbf{x}_{cand}\) satisfies \eqref{eq:assign_one}--\eqref{eq:assign_hum}}
                \State \((\Pi_{cand},T^{tot}_{cand})\leftarrow\Call{Schedule}{\mathbf{x}_{cand},\mathbf{s}}\)
                \If{\((T^{tot},\mathbf{1}^{\mathsf{T}}\mathbf{s})\succ_{\mathrm{lex}}(T^{tot}_{cand},\mathbf{1}^{\mathsf{T}}\mathbf{s})\)}
                    \State \((\mathbf{x},\Pi,T^{tot})\leftarrow(\mathbf{x}_{cand},\Pi_{cand},T^{tot}_{cand})\)
                    \State \(\mathit{improved}\leftarrow\mathrm{true}\); \textbf{break}
                \EndIf
            \EndIf
        \EndFor
    \EndIf
\Until{\(\mathit{improved}=\mathrm{false}\) or \(B_{\mathrm{room}}\) calls of \(\textsc{Schedule}()\)} \label{step:ca-stop}
\State \Return \(\mathbf{x}\)
\end{algorithmic}
\end{algorithm}

In Algorithm~\ref{alg:chamber_assignment}, MIP-CCD(\(\mathbf{s}\)) gives the initial assignment and \(\textsc{Schedule}()\) evaluates it (Steps~\ref{step:ca-exact}--\ref{step:ca-eval}). Chambers are ranked by \(\lambda_c\) from the largest load to the smallest (Step~\ref{step:ca-rank}). We then try \(\textsc{Swap}\) in this heaviest--lightest order (Step~\ref{step:ca-swap}) and accept the first feasible neighbor that improves the lexicographic objective (Step~\ref{step:ca-accept}). If no swap is accepted in that pass, we try \(\textsc{Move}\) in the same order (Step~\ref{step:ca-move}). After an accepted move, \(\lambda_c\) is recomputed and the next pass starts again with \(\textsc{Swap}\). The search stops when a pass yields no improvement or \(B_{\mathrm{room}}\) calls of \(\textsc{Schedule}()\) have been used (Step~\ref{step:ca-stop}).

\subsection{Sample size optimization}
\label{subsec:sample_optimization}
\label{subsec:sample_size_optimization}

For a fixed chamber assignment \(\mathbf{x}\), this stage changes the
sample vector \(\mathbf{s}\).
An extra sample adds one pull-down operation but can let later tests
of that product run in parallel.
Removing a sample does the opposite: pull-down work falls, and so
does the chance of overlap.
Which direction helps is not known in advance, so both are tried
under the same assignment.

The search keeps every product between \(s_{\min}\) and \(s_{\max}\)
samples, with \(s_{\min}=3\) equal to the laboratory planning floor
justified in the problem definition.
The local search starts at this floor and does not go below it.
At the other end, at most eleven tests of a product can run in
parallel, so a twelfth sample cannot create additional concurrent
work.
The upper bound is therefore \(s_{\max}=11\).

The neighborhood changes one product at a time by one or two
samples, provided the new count stays between \(s_{\min}\) and
\(s_{\max}\).
For each product the improving trial with the best tardiness---or,
when tardiness is equal, the smaller sample count---is kept.
After this single-product pass, a Late--Early Rebalance pairs late
products, ordered by decreasing tardiness, with early products,
ordered by decreasing slack, up to \(K_{\mathrm{reb}}\) pairs.
Each pair moves one sample from an early product to a late product
without leaving the bounds.
The paired move is accepted if total tardiness does not increase;
when tardiness is unchanged, the total number of samples must not
increase.
Every trial is decoded by the same earliest-due-date (EDD)
scheduler used in the outer framework.
Algorithm~\ref{alg:sample_ls} states the procedure.

\begin{algorithm}[H]
\caption{SampleSizeLocalSearch}
\label{alg:sample_ls}
\begin{algorithmic}[1]
\State \textbf{Input:} chamber assignment \(\mathbf{x}\), sample vector \(\mathbf{s}\)
\State \textbf{Output:} \((\mathbf{s},\Pi,T^{tot})\)
\State Restrict every \(s_p\) to \([s_{\min},s_{\max}]\) \label{step:ls-init}
\State \((\Pi,T^{tot})\leftarrow\text{Schedule}(\mathbf{x},\mathbf{s})\) \label{step:ls-eval}
\Repeat \label{step:ls-loop}
    \For{each product \(p\in\mathcal{P}\)}
        \State For \(\delta\in\{-2,-1,+1,+2\}\), set \(\mathbf{s}_{cand}\) and \((\Pi_{cand},T^{tot}_{cand})\leftarrow\mathrm{Schedule}(\mathbf{x},\mathbf{s}_{cand})\) \label{step:ls-nbhd}
        \If{\((T^{tot},\mathbf{1}^{\mathsf{T}}\mathbf{s})\succ_{\mathrm{lex}}(T^{tot}_{cand},\mathbf{1}^{\mathsf{T}}\mathbf{s}_{cand})\) for the best trial} \label{step:ls-accept}
            \State \((\mathbf{s},\Pi,T^{tot})\leftarrow(\mathbf{s}_{cand},\Pi_{cand},T^{tot}_{cand})\)
        \EndIf
    \EndFor
    \State \(\mathbf{s}_{reb}\leftarrow\text{LateEarlyRebalance}(\mathbf{s},K_{\mathrm{reb}})\) \label{step:ls-reb}
    \State \((\Pi_{reb},T^{tot}_{reb})\leftarrow\mathrm{Schedule}(\mathbf{x},\mathbf{s}_{reb})\)
    \If{\((T^{tot},\mathbf{1}^{\mathsf{T}}\mathbf{s})\succeq_{\mathrm{lex}}(T^{tot}_{reb},\mathbf{1}^{\mathsf{T}}\mathbf{s}_{reb})\)}
        \State \((\mathbf{s},\Pi,T^{tot})\leftarrow(\mathbf{s}_{reb},\Pi_{reb},T^{tot}_{reb})\)
    \EndIf
\Until{no move was accepted in this pass} \label{step:ls-stop}
\State \Return \((\mathbf{s},\Pi,T^{tot})\)
\end{algorithmic}
\end{algorithm}

Step~\ref{step:ls-init} enforces \(s_{\min}=3\) and \(s_{\max}=11\).
Step~\ref{step:ls-eval} builds the EDD schedule of the current
vector; this is the incumbent of the local search.
Each pass (Step~\ref{step:ls-loop}) first tries the one-product
moves (Step~\ref{step:ls-nbhd}) and accepts a product's best
improving trial (Step~\ref{step:ls-accept}).
Step~\ref{step:ls-reb} then shifts one sample from an early product
to a late product.
Rebalance uses \(\succeq_{\mathrm{lex}}\); the single-product moves use
\(\succ_{\mathrm{lex}}\).
The search stops when a pass accepts nothing
(Step~\ref{step:ls-stop}).

\subsection{Scheduling procedure}
\label{subsec:scheduling_procedure}
\label{subsec:scheduling}

An integral component of the proposed solution approach is to evaluate the quality of a solution. We need to determine the schedule of the tests and the total tardiness of the schedule for a given chamber configuration assignment \(\mathbf{x}\) and a sample-size vector \(\mathbf{s}\). For that purpose we use a constructive scheduling subroutine, called \(\textsc{Schedule}{}\), that assigns tests to stations. Such constructive heuristics are particularly attractive for unrelated parallel-machine problems with eligibility constraints and due-date pressure, where repeated implementation must remain computationally efficient \citep{PerezGonzalezFernandezViagasZamoraGarciaFraminan2019,WangHeKerkhoveVanhoucke2017}.

The main idea in the proposed scheduling procedure is to form a set of tests (called the \emph{candidate test pool}) and update this set at every iteration. This set includes the tests whose precedence relationships are satisfied and that can be scheduled on a station of a compatible chamber whenever a sample is available. At each iteration the procedure rebuilds a candidate pool \(\mathcal{O}^{\mathrm{cand}}\). For each product \(p\), \(\textsc{ReadyTests}(p)\) returns the set of such tests; this set may contain more than one test of the same product, and all of them enter \(\mathcal{O}^{\mathrm{cand}}\). The tests in the pool are then considered in earliest-due-date order. If two tests have the same due date, the one with the earlier feasible start is chosen. If the feasible start times are also equal, the test with the smaller product index, and then the smaller test index, is chosen. The earliest feasible start of a test is the maximum of (i)~the time its precedence constraints are satisfied, (ii)~the time a sample becomes available, and (iii)~the time a station of a compatible chamber becomes available. Only one test is assigned in each iteration. Unselected tests occupy no resource; if they remain ready, they re-enter the pool at the next iteration. The procedure terminates when \(\mathcal{O}^{\mathrm{cand}}=\emptyset\), which for a feasible instance means that all tests have been assigned. The resulting schedule and the total tardiness of the products are returned. We present the steps of the scheduling algorithm in Algorithm~\ref{alg:schedule}.

\begin{algorithm}[H]
\caption{\textsc{Schedule}()}
\label{alg:schedule}
\label{alg:scheduling}
\label{alg:scheduling_procedure}
\label{alg:scheduler_short}
\begin{algorithmic}[1]
\State \textbf{Input:} chamber configuration assignment \(\mathbf{x}\), sample-size vector \(\mathbf{s}\)
\State \textbf{Output:} a feasible schedule \(\Pi\), total tardiness \(T^{tot}\)
\State \(\Pi\leftarrow\emptyset\) \label{step:sch-init}
\Repeat \label{step:sch-loop}
    \State \(\mathcal{O}^{\mathrm{cand}}\leftarrow\emptyset\) \label{step:sch-pool}
    \For{\(p\in\mathcal{P}\)} \label{step:sch-loop1}
        \State \(\mathcal{S}\leftarrow\Call{ReadyTests}{p}\) \label{step:sch-ready}
        \State \(\mathcal{O}^{\mathrm{cand}}\leftarrow\mathcal{O}^{\mathrm{cand}}\cup\mathcal{S}\)
    \EndFor \label{step:sch-loop2}
    \If{\(\mathcal{O}^{\mathrm{cand}}\neq\emptyset\)} \label{step:sch-select}
        \State \(o^{\star}\leftarrow\Call{EarliestDueDate}{\mathcal{O}^{\mathrm{cand}}}\) \label{step:sch-select2}
        \State \(\mathcal{O}^{\mathrm{cand}}\leftarrow\mathcal{O}^{\mathrm{cand}}\setminus\{o^{\star}\}\) \label{step:sch-select3}
        \State \(\Pi\leftarrow\Pi\cup\Call{Assign}{o^{\star}}\) \label{step:sch-assign}
    \EndIf
\Until{\(\mathcal{O}^{\mathrm{cand}}=\emptyset\)} \label{step:sch-stop}
\State Compute \(C_p\), \(T_p\), and \(T^{tot}\) \label{step:sch-obj}
\State \Return \((\Pi,T^{tot})\)
\end{algorithmic}
\end{algorithm}

In Algorithm~\ref{alg:schedule}, we start with an empty schedule (Step~\ref{step:sch-init}). Then, at each iteration we rebuild the candidate test pool: for each product we determine the set of ready tests and add them to \(\mathcal{O}^{\mathrm{cand}}\) (Steps~\ref{step:sch-loop1}--\ref{step:sch-loop2}). \(\textsc{ReadyTests}(p)\) may return more than one test of product \(p\). We then choose a single test from the pool by earliest due date (Steps~\ref{step:sch-select2}--\ref{step:sch-select3}); ties are broken by the earliest feasible start and, if those times are also equal, by the smaller product and test indices. The chosen test is scheduled on a compatible station given the earliest available times of that station and a sample and the time the precedence requirements are fulfilled (Step~\ref{step:sch-assign}). The algorithm terminates when the rebuilt pool is empty (Step~\ref{step:sch-stop}).

\subsection{Variable Neighborhood Search Mechanism}
\label{subsec:vns}
\label{subsec:vns_diversification}

If \(\mathrm{SampleSizeLocalSearch}\) does not improve the incumbent,
Algorithm~\ref{alg:framework} calls \(\mathrm{VNS}\) on the incumbent
chamber assignment and sample vector (Step~\ref{step:fw-vns}).
Following the classical VNS principle of changing neighborhood
intensity after local stagnation \citep{LiaoCheng2007}, and unlike
conventional VNS approaches that perturb sequencing or assignment
decisions, this mechanism perturbs sample-size decisions and therefore
changes the workload structure itself.

We present the main steps of the proposed VNS procedure in
Algorithm~\ref{alg:vns}.
It takes the incumbent solution as the initial solution and perturbs it
with a single destroy ratio \(r\).
In the experiments, \(r=0.20\).
Section~\ref{subsec:ratio} compares this value with \(0.10\) and \(0.30\).
For \(K_{\max}\) iterations, the algorithm visits the two neighborhoods
sequentially (Steps~\ref{VNS:step1}--\ref{VNS:step7}).
One iteration is this pair of neighborhoods.
The downward move and the upward move are not counted as separate
iterations.
In each neighborhood, \(\mathrm{Shake}\) selects a proportion \(r\) of
products and changes their sample sizes.
The first neighborhood (\(l=1\)) decreases the number of samples of
the selected products to \(s_{\min}\), while the second neighborhood
(\(l=2\)) increases the number of samples to \(u_{\max}\)
(Step~\ref{VNS:step2}).
If the current sample count of a product is larger than \(u_{\max}\),
the upward neighborhood does not select that product and leaves the
count unchanged.
The perturbed solution is then improved using the sample-size local
search in Algorithm~\ref{alg:sample_ls}.
The VNS follows a first-improvement policy: as soon as a locally
improved candidate improves the incumbent under the lexicographic
objective, it is returned to the main framework.
If no improvement is found within \(K_{\max}\) iterations, the
incumbent is returned.

\begin{algorithm}[H]
\caption{VNS}
\label{alg:vns}
\label{alg:vns_fixed_destroy}
\label{alg:vns_short}
\begin{algorithmic}[1]
\State \textbf{Input:} chamber assignment \(\mathbf{x}\), incumbent solution \((\mathbf{s}_{inc},\Pi_{inc},T^{tot}_{inc})\)
\State \textbf{Output:} \((\mathbf{s},\Pi,T^{tot})\)
\For{\(k\leftarrow 1\) to \(K_{\max}\)}
    \For{\(l\leftarrow 1\) to \(2\)} \label{VNS:step1}
        \State \(\mathbf{s}_{cand}\leftarrow\mathrm{Shake}(\mathbf{s}_{inc},l,r)\) \label{VNS:step2}
        \State \((\mathbf{s}_{cand},\Pi_{cand},T^{tot}_{cand})\leftarrow\mathrm{SampleSizeLocalSearch}(\mathbf{x},\mathbf{s}_{cand})\)
        \If{\((T^{tot}_{inc},\mathbf{1}^{\mathsf{T}}\mathbf{s}_{inc})\succ_{\mathrm{lex}}(T^{tot}_{cand},\mathbf{1}^{\mathsf{T}}\mathbf{s}_{cand})\)}
            \State \Return \((\mathbf{s}_{cand},\Pi_{cand},T^{tot}_{cand})\)
        \EndIf
    \EndFor \label{VNS:step7}
\EndFor
\State \Return \((\mathbf{s}_{inc},\Pi_{inc},T^{tot}_{inc})\)
\end{algorithmic}
\end{algorithm}

Together, the chamber-assignment, sample-size optimization, scheduling, and VNS components form an iterative framework in which workload generation, resource configuration, and schedule feasibility are updated jointly. The next section evaluates the contribution of these components through computational experiments.

\section{Computational Study}
\label{sec:computational_study}

This section evaluates the computational performance of the proposed Integrated Chamber Setting and Sample Size Optimization (ICS-SSO) approach. We first describe the generated problem instances, the benchmark methods, and the algorithmic settings and performance measures. We then select the VNS perturbation ratio, assess the benefits of product-specific sample allocation and the individual algorithmic components, and examine the performance of the final method under different due-date and voltage-requirement scenarios. All methods were implemented in Java. The experiments were executed on a computer equipped with an Apple M2 chip. Its 8-core CPU consists of four performance cores with a maximum clock rate of 3.49~GHz and four efficiency cores with a maximum clock rate of 2.42~GHz, together with an 8-core GPU.

\subsection{Instance Generation}
\label{subsec:exp_design}

The computational instances are motivated by the industrial testing environment considered in this study. Each instance represents a multi-project testing problem in which products have heterogeneous test requirements, due dates, and voltage requirements. At a given planning point, approximately 50 products are expected to be under test in the laboratory; accordingly, each instance contains 50 products.

Three due-date scenarios are considered to capture differences in deadline tightness and dispersion within an approximately two-month planning horizon. In due-date scenario~1, due dates are relatively concentrated, ranging from 42 to 57 days, with a median of approximately 50 days. Due-date scenario~2 introduces earlier deadlines and greater dispersion, with due dates ranging from 30 to 69 days and a median of approximately 50 days. Due-date scenario~3 is comparatively relaxed while remaining heterogeneous; its due dates range from 34 to 69 days, with a median of approximately 53 days.

The three voltage scenarios differ in the proportion of products requiring a nonstandard supply voltage. Based on the composition observed in industrial practice, approximately 15\%, 30\%, and 45\% of the products require nonstandard voltage in voltage scenarios~1, 2, and 3, respectively. Such requirements restrict the set of chambers eligible to perform certain tests and may create competition for voltage-capable chambers. The implications for total tardiness, however, also depend on the remaining characteristics of the instance and the decisions made by the solution method.

Ten test-requirement scenarios are evaluated for each combination of the three due-date and three voltage scenarios. Thus, the experiments comprise \(10\times3\times3=90\) instances. For every product, refrigerant charge determination, pull-down, the performance test at \(25^\circ\mathrm{C}\), and the three consumer-usage tests are required. Energy-efficiency tests, the remaining climatic-class performance tests, and the optional freezing-capacity and temperature-rise tests are selected according to the rules described in Table~\ref{tab:test_attributes}. The resulting instance set enables comparisons across heterogeneous test portfolios, deadline patterns, and voltage requirements.

\subsection{Benchmark Methods}
\label{subsec:methods}

We compare ICS-SSO with two uniform sample-size policies and two reduced variants of the proposed method. All approaches use the same test-scheduling procedure and respect the same feasibility requirements. The differences concern whether sample counts are fixed or optimized and whether feedback and diversification are enabled.

The first benchmark, \emph{FS(3)}, represents the current industrial policy of assigning three samples to every product. Three samples are also the minimum accepted in the considered setting. Sample counts remain fixed throughout the procedure; chamber temperature and humidity settings are established once using the chamber-assignment procedure described in Section~\ref{subsec:chamber_assignment}. No sample-size optimization is performed.

The second benchmark, \emph{FS(6)}, assigns six samples to each product, yielding 300 samples for the 50 products in an instance. This benchmark tests whether a uniform increase in sample availability is sufficient to obtain schedules comparable to those of the proposed optimization approach. As in FS(3), the chamber settings are established once and the sample counts are not subsequently modified. FS(6) is used as a uniform-allocation reference; it does not imply that ICS-SSO uses exactly 300 samples in every instance.

The complete method, \emph{ICS-SSO}, jointly explores sample-size decisions and chamber settings through the feedback mechanism described earlier in the paper, together with VNS-based diversification. To assess the feedback mechanism, we consider \emph{CS-SSO-noFM} (Chamber Setting and Sample Size Optimization without Feedback Mechanism). This variant determines chamber settings once and optimizes sample sizes without revisiting those settings. It corresponds to the \emph{no outer} configuration in the computational experiments.

To assess diversification, we consider \emph{ICS-SSO-noVNS} (Integrated Chamber Setting and Sample Size Optimization without Diversification). This variant retains the interaction between chamber-setting and sample-size decisions, but removes the VNS perturbation phase. It therefore relies on feedback and local improvement without diversification of the sample-size vector. These two reduced variants are used separately to identify the contribution associated with each component.

\subsection{Parameter Settings and Performance Measures}
\label{subsec:parameter_settings}

All methods use common parameter values for the components they share. For ICS-SSO, the search begins with three samples per product, consistent with current industrial practice. The sample count of each product is bounded by \(s_{\min}=3\) and \(s_{\max}=11\); the latter reflects the maximum number of tests of a product that can be performed in parallel in the considered setting. A local sample-size move increases or decreases the sample count of one product by one or two units, subject to these bounds.

In the VNS phase, a proportion \(R_{\mathrm{VNS}}\) of the products is selected for perturbation. A downward perturbation resets the selected sample counts to \(s_{\min}\), whereas an upward perturbation resets them to \(u_{\max}=8\); local improvement follows the perturbation. The iteration counter is defined at the VNS iteration level, rather than counting the upward and downward perturbation directions as separate iterations. The maximum VNS iteration budget is fixed at \(K_{\max}=100\). Three fixed perturbation ratios, 0.10, 0.20, and 0.30, are examined in Section~\ref{subsubsec:vns_ratio}; the selected ratio for the final method is 0.20. No switching between perturbation ratios is used within a run.

The lexicographic objective gives priority to minimizing total tardiness. If two solutions have identical total tardiness, the solution with fewer samples is preferred. The Late--Early Rebalance procedure respects this ordering: a move is accepted only if total tardiness does not increase and, when tardiness remains unchanged, total sample usage does not increase. The fixed settings used in the experiments are reported in Table~\ref{tab:param_settings_final}.

\begin{table}[H]
\centering
\caption{Parameter settings used in the computational study.}
\label{tab:param_settings_final}
\begin{tabular}{lll}
\hline
\textbf{Parameter} & \textbf{Description} & \textbf{Value} \\
\hline
\(s_p^{(0)}\) & Initial sample count per product & 3 \\
\(s_{\min}\) & Minimum sample count & 3 \\
\(u_{\max}\) & Upper sample count in a VNS upward move & 8 \\
\(s_{\max}\) & Maximum sample count & 11 \\
\(\Delta_p\) & Sample-count moves & \(\{-2,-1,+1,+2\}\) \\
\(R_{\mathrm{VNS}}\) & Selected VNS perturbation ratio & 0.20 \\
\(I_{\max}\) & Maximum outer-loop iterations & 100 \\
\(K_{\mathrm{reb}}\) & Maximum Late--Early Rebalance pairs & 10 \\
\(K_{\max}\) & Maximum VNS iterations & 100 \\
\(B_{\mathrm{room}}\) & Chamber-assignment local-search budget & 200 \\
\(\rho\) & Random seed for VNS product selection & 42 \\
\hline
\end{tabular}
\end{table}

The computational results are evaluated using four measures. \emph{Average tardiness} is the mean total tardiness over the instances in the corresponding comparison. \emph{Average gap} and \emph{gap standard deviation} summarize the instance-level percentage deviations from the best total-tardiness value within the comparison set. For instance \(i\) and method \(m\), let \(T_{i,m}\) denote the obtained total tardiness and let \(T_i^{\mathrm{best}}\) denote the smallest total tardiness among the alternatives compared for that instance. The percentage gap is
\begin{equation}
\label{eq:gap}
\mathrm{Gap}_{i,m}(\%) = \frac{T_{i,m}-T_i^{\mathrm{best}}}{T_i^{\mathrm{best}}}\times 100.
\end{equation}
The average gap is the arithmetic mean of the instance-level gap values, and the gap standard deviation is their sample standard deviation. These statistics are therefore \emph{not} obtained by calculating the percentage difference between overall mean tardiness values. Finally, \emph{best count} is the number of instances for which a method attains the lowest tardiness in its comparison set. Ties are counted for every method attaining that value, so the sum of best counts may exceed the number of instances. Because gaps are measured relative to the best method within a given comparison, zero gaps and zero gap dispersion for a consistently best-performing method follow from the definition; they do not imply that its total tardiness is constant across instances.

\subsection{Computational Results}
\label{subsec:overall_results}

\subsubsection{Selection of the VNS Perturbation Ratio}
\label{subsubsec:vns_ratio}
\label{subsec:ratio}

We first investigate the sensitivity of ICS-SSO to the VNS perturbation ratio. The algorithm is evaluated with the three fixed ratios \(R_{\mathrm{VNS}}=0.10\), 0.20, and 0.30 on all 90 instances. The iteration budget is held at \(K_{\max}=100\), and the remaining parameters are unchanged. For this comparison, the best tardiness on each instance is the smallest of the values obtained with the three ratios. Table~\ref{tab:vns_ratio_comparison} presents the resulting statistics.

\begin{table}[H]
\centering
\caption{Comparison of VNS perturbation ratios over 90 instances.}
\label{tab:vns_ratio_comparison}
\footnotesize
\setlength{\tabcolsep}{4pt}
\begin{tabular}{lrrrr}
\hline
\textbf{Ratio} & \textbf{Avg. Tardiness} & \textbf{Avg. Gap (\%)} & \textbf{Gap Std. Dev. (\%)} & \textbf{Best Count} \\
\hline
0.10 & 203.62 & 27.68 & 43.37 & 30 \\
0.20 & 189.64 & 12.83 & 21.91 & 42 \\
0.30 & 201.49 & 21.46 & 27.83 & 20 \\
\hline
\end{tabular}
\end{table}

The intermediate perturbation ratio offers the most favorable overall performance among the settings considered, not only in terms of mean solution quality but also in terms of the dispersion of relative gaps. This suggests that a moderate perturbation is better suited to the search than either a more conservative or a more disruptive change to the current sample allocation. A small perturbation may leave the search close to its current local solution, whereas a large perturbation may discard too much of a promising allocation. These are plausible explanations of the observed pattern rather than independently established mechanisms. Since no single ratio dominates every instance, we select \(R_{\mathrm{VNS}}=0.20\) on the basis of its aggregate performance and use it throughout the remaining experiments.

\subsubsection{Comparison with Fixed Sample-Size Policies}
\label{subsubsec:fixed_sample}

We next compare the final ICS-SSO configuration with the two uniform sample-size policies. FS(3) reflects the existing industrial rule, while FS(6) examines the effect of providing twice as many samples uniformly to each product. The instance-wise best tardiness for this comparison is determined from FS(3), FS(6), and ICS-SSO. Table~\ref{tab:sample_policy_comparison} reports the results over all 90 instances.

\begin{table}[H]
\centering
\caption{Comparison of ICS-SSO with uniform sample-size policies over 90 instances.}
\label{tab:sample_policy_comparison}
\footnotesize
\setlength{\tabcolsep}{4pt}
\begin{tabular}{lrrrr}
\hline
\textbf{Method} & \textbf{Avg. Tardiness} & \textbf{Avg. Gap (\%)} & \textbf{Gap Std. Dev. (\%)} & \textbf{Best Count} \\
\hline
FS(3) & 721.01 & 360.32 & 198.59 & 0 \\
FS(6) & 472.83 & 196.95 & 121.44 & 0 \\
ICS-SSO & 189.64 & 0.00 & 0.00 & 90 \\
\hline
\end{tabular}
\end{table}

Increasing the number of samples under a uniform policy improves scheduling performance, indicating that additional opportunities for parallel testing can alleviate some delays. Nevertheless, uniformly increasing sample counts does not close the performance gap with ICS-SSO. The results suggest that the value of additional samples depends on where they are allocated: products differ in their due dates, test sequences, and demands for compatible chamber environments, so providing the same number of samples to every product does not necessarily address the most consequential scheduling bottlenecks.

This comparison establishes the limitations of uniform sample-size policies, but it does not isolate sample allocation as the sole source of the gain, because ICS-SSO also includes chamber-setting feedback and VNS diversification. The following ablation experiments investigate these additional components. Moreover, the present results do not support a claim of lower total sample consumption by ICS-SSO, since its realized sample totals are not reported.

\subsubsection{Contribution of the Feedback Mechanism}
\label{subsubsec:feedback}

To quantify the contribution of the feedback mechanism, the final ICS-SSO method is compared with CS-SSO-noFM. In this variant, chamber settings are fixed after the initial assignment and are not reconsidered when the sample-size decisions change. The comparison therefore evaluates the value of revisiting chamber conditions as sample allocation evolves. Table~\ref{tab:feedback_comparison} summarizes the results.

\begin{table}[H]
\centering
\caption{Effect of the feedback mechanism over 90 instances.}
\label{tab:feedback_comparison}
\footnotesize
\setlength{\tabcolsep}{4pt}
\begin{tabular}{lrrrr}
\hline
\textbf{Method} & \textbf{Avg. Tardiness} & \textbf{Avg. Gap (\%)} & \textbf{Gap Std. Dev. (\%)} & \textbf{Best Count} \\
\hline
CS-SSO-noFM & 394.36 & 139.16 & 90.22 & 0 \\
ICS-SSO & 189.64 & 0.00 & 0.00 & 90 \\
\hline
\end{tabular}
\end{table}

The deterioration observed when feedback is removed indicates that treating chamber settings as a one-time decision can be restrictive, even when sample sizes are subsequently optimized. Sample-size changes modify the workload assigned to different environmental conditions, including workload associated with pull-down tests. A chamber configuration that is appropriate for an earlier sample allocation may therefore become less suitable as the search progresses. Revisiting chamber settings allows these interdependent decisions to be adjusted together.

The consistent advantage of the full method across the test set supports the inclusion of feedback as part of the integrated search. However, the comparison evaluates the resulting scheduling performance; without chamber-configuration histories, it does not identify which individual setting changes account for the improvement.

\subsubsection{Contribution of VNS Diversification}
\label{subsubsec:vns_contribution}

The role of VNS diversification is examined by comparing ICS-SSO with ICS-SSO-noVNS. Both methods preserve feedback between chamber-setting and sample-size optimization, but ICS-SSO-noVNS omits the perturbation phase after local improvement. This comparison isolates the performance contribution associated with diversification of sample-size decisions. The corresponding results are given in Table~\ref{tab:diversification_comparison}.

\begin{table}[H]
\centering
\caption{Effect of VNS diversification over 90 instances.}
\label{tab:diversification_comparison}
\footnotesize
\setlength{\tabcolsep}{4pt}
\begin{tabular}{lrrrr}
\hline
\textbf{Method} & \textbf{Avg. Tardiness} & \textbf{Avg. Gap (\%)} & \textbf{Gap Std. Dev. (\%)} & \textbf{Best Count} \\
\hline
ICS-SSO-noVNS & 444.91 & 177.67 & 113.53 & 0 \\
ICS-SSO & 189.64 & 0.00 & 0.00 & 90 \\
\hline
\end{tabular}
\end{table}

Removing VNS diversification leads to a pronounced deterioration in solution quality despite retaining the feedback mechanism. This suggests that repeatedly optimizing sample sizes and chamber settings from the current solution may not provide sufficient exploration of the feasible search space. By perturbing the sample-size vector and then restarting local improvement, VNS introduces alternative allocation patterns that the unperturbed search may not reach.

The dispersion of the relative gaps also indicates that the extent of the performance loss without VNS varies among instances. The experiment therefore supports diversification as a useful complement to feedback and local improvement, rather than implying that it provides a fixed amount of improvement in every case. The effects of removing feedback and removing VNS are separate ablations and should not be added as if they were independent percentage contributions.

\subsubsection{Performance Across Due-Date Scenarios}
\label{subsubsec:due_effect}

We next evaluate how the final ICS-SSO method performs across the three due-date scenarios. Each scenario comprises 30 instances, corresponding to the ten test scenarios and three voltage scenarios. In this analysis, the comparisons are restricted to the two fixed-sample policies and ICS-SSO, as shown in Table~\ref{tab:due_scenario_comparison}. For each group, the gaps are computed at the instance level relative to the smallest tardiness among these three methods.

\begin{table}[H]
\centering
\caption{Performance under different due-date scenarios (30 instances per scenario).}
\label{tab:due_scenario_comparison}
\footnotesize
\setlength{\tabcolsep}{3.2pt}
\begin{tabular}{llrrrr}
\hline
\textbf{Scenario} & \textbf{Method} & \textbf{Avg. Tardiness} & \textbf{Avg. Gap (\%)} & \textbf{Gap Std. Dev. (\%)} & \textbf{Best Count} \\
\hline
Due 1 & FS(3) & 744.17 & 292.40 & 88.75 & 0 \\
 & FS(6) & 480.20 & 154.16 & 61.02 & 0 \\
 & ICS-SSO & 199.97 & 0.00 & 0.00 & 30 \\
\hline
Due 2 & FS(3) & 803.37 & 211.10 & 64.66 & 0 \\
 & FS(6) & 559.67 & 117.19 & 46.76 & 0 \\
 & ICS-SSO & 272.27 & 0.00 & 0.00 & 30 \\
\hline
Due 3 & FS(3) & 615.50 & 577.46 & 179.98 & 0 \\
 & FS(6) & 378.63 & 319.49 & 123.87 & 0 \\
 & ICS-SSO & 96.70 & 0.00 & 0.00 & 30 \\
\hline
\end{tabular}
\end{table}

The due-date scenarios exhibit a clear difference in scheduling difficulty. The scenario with earlier and more dispersed deadlines yields the highest average tardiness, whereas the comparatively relaxed scenario yields the lowest. The same ordering appears under both fixed-sample policies and ICS-SSO, suggesting that the deadline structure affects the underlying scheduling problem rather than only one particular solution procedure.

ICS-SSO remains the best-performing approach within every due-date group, including the more demanding conditions. Its relative advantage is especially pronounced under the relaxed scenario, where the optimized solutions have relatively small tardiness values. Because the percentage-gap measure uses those values as its denominator, large relative gaps need not correspond to equally large absolute differences. Considering absolute tardiness alongside the gaps avoids overstating the effect of due-date conditions on the size of the improvement.

\subsubsection{Performance Across Voltage Scenarios}
\label{subsubsec:voltage_effect}

Finally, we investigate the performance of ICS-SSO for different proportions of products requiring nonstandard voltage. Voltage scenarios~1, 2, and 3 correspond to approximately 15\%, 30\%, and 45\% of the products, respectively. Each voltage group contains 30 instances, covering the ten test scenarios and three due-date scenarios. Table~\ref{tab:voltage_scenario_comparison} compares the two fixed-sample policies with the final ICS-SSO method and reports gaps relative to the instance-wise best of these three methods.

\begin{table}[H]
\centering
\caption{Performance under different voltage scenarios (30 instances per scenario).}
\label{tab:voltage_scenario_comparison}
\footnotesize
\setlength{\tabcolsep}{3.2pt}
\begin{tabular}{llrrrr}
\hline
\textbf{Scenario} & \textbf{Method} & \textbf{Avg. Tardiness} & \textbf{Avg. Gap (\%)} & \textbf{Gap Std. Dev. (\%)} & \textbf{Best Count} \\
\hline
Volt 1 (15\%) & FS(3) & 716.07 & 337.92 & 175.48 & 0 \\
 & FS(6) & 462.77 & 177.01 & 101.71 & 0 \\
 & ICS-SSO & 199.50 & 0.00 & 0.00 & 30 \\
\hline
Volt 2 (30\%) & FS(3) & 720.57 & 389.25 & 265.81 & 0 \\
 & FS(6) & 470.20 & 214.20 & 164.91 & 0 \\
 & ICS-SSO & 189.97 & 0.00 & 0.00 & 30 \\
\hline
Volt 3 (45\%) & FS(3) & 726.40 & 353.79 & 134.76 & 0 \\
 & FS(6) & 485.53 & 199.64 & 83.64 & 0 \\
 & ICS-SSO & 179.47 & 0.00 & 0.00 & 30 \\
\hline
\end{tabular}
\end{table}

The fixed-sample policies show a gradual increase in average tardiness as the proportion of products requiring nonstandard voltage rises. This pattern is consistent with the additional competition for voltage-capable chambers that can arise when chamber eligibility becomes more restrictive. The full ICS-SSO method, however, maintains its advantage across all voltage groups and exhibits a different trend in group-average tardiness.

The contrast highlights the relevance of coordinated sample-allocation and chamber-setting decisions when the test portfolio includes voltage-dependent eligibility restrictions. It should not, however, be interpreted as evidence that a higher nonstandard-voltage share makes the problem intrinsically easier or necessarily benefits ICS-SSO. The reported averages describe the generated groups, and changes in the product mix and other instance characteristics can also influence the outcome. What the experiment establishes is that the advantage over uniform sample policies persists across the tested voltage requirements.

\section{Conclusion}
\label{sec:conclusion}

This study introduced the Test Chamber Scheduling with Sample Optimization Problem (TCSSOP), motivated by refrigerator R\&D testing operations. The problem combines heterogeneous chamber resources, temperature-humidity requirements, voltage and humidity eligibility, stage-based precedence logic, due dates, and sample-size-dependent workload generation. Unlike classical scheduling problems with fixed job sets, TCSSOP treats project-level sample counts as decision variables that affect both upstream pull-down workload and downstream parallelization opportunities.

To address this problem, an iterative framework was proposed. For a given sample vector, chamber environments are assigned according to workload distribution and chamber capabilities, then refined by local search. Given chamber settings, sample sizes are optimized through bounded local search and evaluated by an EDD-based constructive scheduler with explicit feasibility validation. Because sample decisions alter environment-level workload, chamber assignments are updated iteratively. When local improvement stagnates, a VNS diversification mechanism perturbs sample-size decisions to explore alternative workload structures rather than perturbing sequencing or assignment alone.

Computational experiments on 90 benchmark scenarios demonstrate that optimized sample allocation yields substantial tardiness reductions relative to uniform industrial policies. The selected configuration uses a single VNS destroy ratio of \(0.20\) with \(K_{\max}=100\), where one iteration applies both the downward and the upward neighborhood. This configuration obtains the lowest total tardiness in all 90 instances. Relative to the fixed three-sample baseline FS(3), the average deviation is \(360.32\%\). The uniform six-sample policy FS(6) reduces tardiness but still has an average deviation of \(196.95\%\). Removing chamber-assignment feedback (CS-SSO-noFM) or VNS diversification (ICS-SSO-noVNS) leaves average deviations of \(139.16\%\) and \(177.67\%\), respectively. Sample-size optimization, iterative chamber reconfiguration, and VNS diversification are complementary.

From a managerial perspective, prototype sample size should be treated as a tactical planning decision rather than a fixed convention applied uniformly across projects. Raising every product from three samples to six lowers tardiness, but the uniform six-sample rule remains substantially inferior to selective allocation. Additional samples should therefore be directed to projects for which downstream parallelism and due-date pressure justify the upstream pull-down cost, while chamber configurations should be reassessed whenever sample decisions change. The lexicographic objective also discourages unnecessary prototype consumption by preferring lower total sample usage among tardiness-equivalent solutions. Finally, the results emphasize the importance of feasibility validation, since schedules that violate chamber compatibility, station capacity, sample exclusivity, or stage-based precedence logic are not operationally implementable.

This study has several limitations. The benchmark instances are synthetically generated to preserve confidentiality while reflecting industrial structure; validation on anonymized real project portfolios remains desirable. Chamber environments are treated as fixed over the planning horizon, and test durations are deterministic. Setup times for chamber reconfiguration are not modeled explicitly, and the reported gaps measure relative performance among the evaluated methods rather than optimality gaps against an exact benchmark.

Future research can extend TCSSOP in several directions. Stochastic test durations, test failures, retesting requirements, and chamber breakdowns can be incorporated to improve robustness. Chamber setup times can be modeled explicitly if chambers are allowed to change environments over time. Multi-objective extensions may consider energy consumption, labor availability, prototype cost, and testing cost together with tardiness and sample usage. Exact or matheuristic methods on smaller instances could provide optimality gaps for the proposed heuristic. Finally, rolling-horizon versions of the problem could support dynamic project arrivals in industrial R\&D environments.

\bibliographystyle{plainnat}
\bibliography{References}

\end{document}
